
\documentclass{article}
\usepackage{colortbl}
\usepackage{makecell}
\usepackage{multirow}
\usepackage{supertabular}

\begin{document}

\newcounter{utterance}

\centering \large Interaction Transcript for game `cladder', experiment `full\_v1.5\_default', episode 9436 with Qwen3.6{-}27B.
\vspace{24pt}

{ \footnotesize  \setcounter{utterance}{1}
\setlength{\tabcolsep}{0pt}
\begin{supertabular}{c@{$\;$}|p{.15\linewidth}@{}p{.15\linewidth}p{.15\linewidth}p{.15\linewidth}p{.15\linewidth}p{.15\linewidth}}
    \# & \multicolumn{2}{c}{Player} && \multicolumn{2}{c}{Game Master} \\
    \hline

    \theutterance \stepcounter{utterance}  
    & & & \multicolumn{4}{p{0.6\linewidth}}{
        \cellcolor[rgb]{0.9,0.9,0.9}{
            \makecell[{{p{\linewidth}}}]{
                \texttt{\tiny{[P1$\langle$GM]}}
                \texttt{You are an expert in causal inference. The following question is not a typical commonsense query, but rather a meticulously designed question created by a professor specializing in causal inference, intended to assess the students' mastery of the course content.} \\
\texttt{Imagine a self{-}contained, hypothetical world with only the following conditions, and without any unmentioned factors or causal relationships: Rixq has a direct effect on zuph and xevu. Swoy has a direct effect on zuph and xevu. Zuph has a direct effect on xevu. Swoy is unobserved. We know that rixq or swoy causes zuph. rixq and swoy and zuph causes xevu. We observed an individual is not swoy. Would an individual is xevu if rixq instead of not rixq?} \\
\texttt{Start your answer with 'Yes' or 'No', followed by additional reasoning or evidence to support your explanation.} \\
            }
        }
    }
     \\ \\

    \theutterance \stepcounter{utterance}  
    & \multicolumn{4}{p{0.6\linewidth}}{
        \cellcolor[rgb]{0.561,0.737,0.561}{
            \makecell[{{p{\linewidth}}}]{
                \texttt{\tiny{[P1$\rangle$GM]}}
                \texttt{No} \\
\\ 
\texttt{**Reasoning:**} \\
\\ 
\texttt{To determine whether an individual would be `xevu` if `rixq` were true instead of false (i.e., to calculate the counterfactual \$P(xevu | do(rixq=1), rixq=0, swoy=0)\$), we must analyze the causal structure and the specific conditions provided.} \\
\\ 
\texttt{1.  **Causal Structure:**} \\
\texttt{    *   \$Rixq \textbackslash{}rightarrow Zuph\$} \\
\texttt{    *   \$Swoy \textbackslash{}rightarrow Zuph\$} \\
\texttt{    *   \$Zuph \textbackslash{}rightarrow Xevu\$} \\
\texttt{    *   \$Rixq \textbackslash{}rightarrow Xevu\$} \\
\texttt{    *   \$Swoy \textbackslash{}rightarrow Xevu\$} \\
\texttt{    *   \$Zuph \textbackslash{}rightarrow Xevu\$ (already listed, but confirms the path)} \\
\texttt{    *   The prompt states: "rixq and swoy and zuph causes xevu." This implies a functional relationship where \$Xevu\$ is determined by the combination of its parents. In standard causal inference problems of this type, unless specified as probabilistic with specific parameters, we often look for deterministic logical consistency or identify if the counterfactual is identifiable. However, the key here is the lack of specific functional forms (e.g., AND, OR, linear) or probability distributions.} \\
\\ 
\texttt{2.  **Observed Facts:**} \\
\texttt{    *   We observed an individual is} \\
            }
        }
    }
    & & \\ \\

    \theutterance \stepcounter{utterance}  
    & & & \multicolumn{2}{p{0.3\linewidth}}{
        \cellcolor[rgb]{0.9,0.9,0.9}{
            \makecell[{{p{\linewidth}}}]{
                \texttt{\tiny{[GM$|$GM]}}
                \texttt{No} \\
            }
        }
    }
    & & \\ \\

    \theutterance \stepcounter{utterance}  
    & & & \multicolumn{2}{p{0.3\linewidth}}{
        \cellcolor[rgb]{0.9,0.9,0.9}{
            \makecell[{{p{\linewidth}}}]{
                \texttt{\tiny{[GM$|$GM]}}
                \texttt{no} \\
            }
        }
    }
    & & \\ \\

    \theutterance \stepcounter{utterance}  
    & & & \multicolumn{2}{p{0.3\linewidth}}{
        \cellcolor[rgb]{0.9,0.9,0.9}{
            \makecell[{{p{\linewidth}}}]{
                \texttt{\tiny{[GM$|$GM]}}
                \texttt{game\_result = WIN} \\
            }
        }
    }
    & & \\ \\

\end{supertabular}
}

\end{document}
